\documentclass[10pt,a4paper]{amsart}
\usepackage[margin=19mm]{geometry}
\usepackage{amsmath,amssymb,mathtools}
\usepackage[scr=rsfs,cal=euler]{mathalfa}
\usepackage{xfrac}
\usepackage[unicode,hidelinks,bookmarks=false]{hyperref}
\newcommand{\ie}{i.e.\ }
\newcommand{\cf}{cf.\ }
\newcommand{\resp}{respectively\ }
\newcommand{\Subsection}{\subsection}
\providecommand{\nobreakdash}{\nobreak\hbox{-}}
\setlength{\emergencystretch}{2em}
\multlinegap0pt
\usepackage{bm}
\usepackage{bbm}
\usepackage{dsfont}
\usepackage{xspace}
\usepackage{cancel}
\usepackage{mathtools}
\usepackage{amsthm}
\makeatletter
\@ifundefined{defi}{\newtheorem{defi}{Defi}}{}
\@ifundefined{prop}{\newtheorem{prop}[defi]{Proposition}}{}
\makeatother
\def\PP{\mathcal{P}}
\def\R{\mathbb{R}}
\def\SS{\mathcal{S}}
\title[Uniqueness problem for Prandtl spirals]{Uniqueness problem for Prandtl spirals}
\author{Tomasz Cie\'slak, Piotr Kokocki, Przemys{\l}aw Kosewski}
\date{Source: arXiv:2508.03554v1 (CC BY 4.0)}
\begin{document}
\maketitle

\noindent\emph{Source and licence.} Uniqueness problem for Prandtl spirals. Version arXiv:2508.03554v1 (CC BY 4.0), distributed under the Creative Commons Attribution 4.0 International licence, as recorded in the arXiv licence metadata for this exact version and shown on the abstract page https://arxiv.org/abs/2508.03554v1. Full rights notice: licenses/arxiv-2508.03554-CC-BY-4.0.txt.\par\smallskip

\noindent\textit{Editorial scope.} Counted unit: the Prop whose source label is \texttt{tresc} in the article (source0.tex lines 438--444, source label \texttt{tresc}), delivered together with the article's own complete original proof of it (source0.tex lines 445--518, one \texttt{proof} environment of 74 lines). No mathematics is written, repaired or re-derived: statement and proof are the article's own text line for line. Required context retained uncounted: holojedno (lines 324--330); eq-36 (lines 424--427); eq-39 (lines 424--427). Every definition, display and label of those lines is retained unchanged. Omitted from this package and not redistributed: 1--323, 331--423, 428--437, 519--1191. Nothing inside a retained excerpt is omitted or altered: every retained body is delivered exactly as the article's lines are written, so no omission receipt of any kind accompanies this package. Citations: the retained excerpts carry no citation command at all, so no bibliography entry of the article is transcribed and no reference is redistributed. Reference adaptation: none was needed and none was made; every reference inside the delivered unit resolves inside it, so no unresolved reference or citation is produced. Printed slips kept literal: no printed word, symbol, hypothesis or number of the counted statement or of its proof was altered. Layout and class: the article's own class is \texttt{amsart} with packages including amssymb, amsmath, mathrsfs, amsfonts, amsthm, graphicx, subcaption, hyperref, verbatim, pdfsync, lmodern, babel; the wrapper is the lane's pinned \texttt{amsart} editorial wrapper at 10pt on A4, which restates the theorem-like environments the retained text needs and adds the article's own macro definitions that the retained text needs (\texttt{\textbackslash PP}, \texttt{\textbackslash R}, \texttt{\textbackslash SS}), with extra wrapper packages (bm, bbm, dsfont, xspace, cancel, mathtools). The delivered numbering is the unit's own. Assets: no retained span contains a figure environment, a graphics-inclusion command, a PDF-page inclusion, a table or a TikZ picture; no image file is copied into this package, the package declares no asset dependency and no image file is redistributed with it. Licence: version arXiv:2508.03554v1 of this article is distributed under the Creative Commons Attribution 4.0 International licence, as recorded in the arXiv licence metadata for this exact version and shown on the abstract page https://arxiv.org/abs/2508.03554v1. A full rights notice is delivered at licenses/arxiv-2508.03554-CC-BY-4.0.txt. Characters: every retained excerpt is pure ASCII, so the delivered engine, XeLaTeX with its default Latin Modern text font, reports no missing character.
\begin{prop}\label{holojedno}
Suppose $h$ is a holomorphic function defined on the open vertical strip
\[
S_{x_{0},x_{1}} := \{x + it \,:\, x \in (x_0, x_1), \, t \in \R\},
\]
where $x_0 < x_1$, and assume that $f$ extends continuously to the boundary of $S$. If $h(x_1 + it) = 0$ for all $t \in \R$, then $h \equiv 0$ on the entire strip $S_{x_{0},x_{1}}$.
\end{prop}

\begin{align}\label{eq-36}
& \tilde{w}_{1}(iy) = \mu e^{2ay}, \quad \tilde{w}_{1}\left(-\frac{2a\pi}{1+a^2} + iy\right) = \mu e^{2a(y - \frac{2\pi}{1+a^2})}, \\ \label{eq-39}
& \tilde{w}_{2}\left(-\frac{2a\pi}{1+a^2} + i\left(y + \frac{2\pi}{1+a^2}\right)\right) - \tilde{w}_{2}(iy) = 2ag e^{2ay}
\end{align}

\begin{prop}\label{tresc}
Let $\tilde{w}^j = \tilde{w}^{j}_{1}+i\tilde{w}^{j}_{2}$, for $j = 1, 2$, be two antiholomorphic velocities in the strip $\SS$ satisfying the boundary conditions \eqref{eq-36} and the non-standard condition \eqref{eq-39}. Furthermore, assume that for any $j=1,2$ and $-2\pi a/(1+a^2)<x<0$ we have the pointwise limit
\begin{equation}\label{obietnica}
\tilde{w}^j(x+iy)\rightarrow 0\quad \text{as} \ \ y\to -\infty.
\end{equation}
Then $\tilde{w}^1 = \tilde{w}^2$ on $\SS$.
\end{prop}

\begin{proof}
Let us consider the mappings 
\begin{align*}
\PP_{\pm}(z) = -z^{*} - \frac{2a\pi}{1+a^2} \pm \frac{2\pi i}{1+a^2}, \quad z\in\SS, 
\end{align*}
which are well-defined transformations of the strip $\SS$ onto itself. More precisely, each $\PP_{\pm}$ is constructed as the composition of a reflection with respect to the symmetry axis of the strip $\{\mathrm{Re}\,z = - \tfrac{a\pi}{1+a^2}\}$ followed by a vertical shift by the vector $\pm\tfrac{2\pi i}{1+a^2}$. Moreover, the following equality holds 
\begin{equation}\label{eq-p1}
\PP_{+}(\PP_{-}(z)) = z, \quad z\in\SS
\end{equation}
and 
\begin{equation}\label{eq-p2}
\PP_{-}^{n}(z) = z - \frac{2\pi n i}{1+a^2} \quad \text{ for } z\in\SS\text{ and even } n\ge 1.
\end{equation}
Let us define the auxiliary functions $$g^j(z) := \tilde{w}_{1}^j(\PP_{+}(z)) + i\tilde{w}_{2}^j(\PP_{+}(z)), \quad z\in\SS, \ j=1,2.$$ 
Since the function $(w^{j})^{*}$ is holomorphic, the maps $g^j$ also have this property, as they satisfy the following Cauchy-Riemann equations 
\begin{equation*}
(\mathrm{Re}\, g^{j})_{x}(z) = -\tilde{w}_{1,x}^j(\PP_{+}(z)) = \tilde{w}_{2,y}^j(\PP_{+}(z)) = (\mathrm{Im}\, g^{j})_{y}(z)
\end{equation*}
and
\begin{equation*}
(\mathrm{Re}\, g^{j})_{y}(z) = \tilde{w}_{1,y}^j(\PP_{+}(z)) = \tilde{w}_{2,x}^j(\PP_{+}(z)) = -(\mathrm{Im}\, g^{j})_{x}(z).
\end{equation*}
Then we can define the holomorphic function on the strip $\SS$ as follows
\begin{equation*}
\begin{aligned}
h^j(z) & := (w^j(z))^{*} + g^j(z) \\
& = \tilde{w}_{1}^j(z) + \tilde{w}_{1}^j(\PP_{+}(z)) - i \left(\tilde{w}_{2}^j(z) - \tilde{w}_{2}^j(\PP_{+}(z))\right).
\end{aligned}
\end{equation*}
By \eqref{eq-36} and \eqref{eq-39}, the following boundary conditions are satisfied 
\begin{equation}
\begin{aligned}\label{eq-50}
h^j(iy) &= \tilde{w}_{1}^j(iy) + \tilde{w}_{1}^j\left(-\frac{2a\pi}{1+a^2}+ i\left(y + \frac{2\pi}{1+a^2}\right)\right) \\
&\quad - i \left(\tilde{w}_{2}^j(iy) - \tilde{w}_{2}^j\left(-\frac{2a\pi}{1+a^2}+ iy + \frac{2\pi i}{1+a^2}\right)\right) \\
&= 2\mu e^{2ay} +i 2a g e^{2ay} \quad \text{for} \ \ y\in\R.
\end{aligned}
\end{equation}
Let us now define the function
\begin{equation*}
h(z) := h^1(z) - h^2(z), \quad z\in\SS,
\end{equation*}
which is holomorphic on $\SS$ and continuous up to the boundary $\partial\SS$. Moreover, by \eqref{eq-50}, it is known that $h(iy) = 0$ for $y\in\R$. By applying the Schwarz reflection principle (see Proposition \ref{holojedno}), we conclude that $h^1(z) - h^2(z) = h(z)=0$ for $z\in\SS$. 

To prove the uniqueness of the velocity field, let us note that
\begin{equation}\label{formula-1}
\tilde{w}_{2}^j(\PP_{+}(z)) - \tilde{w}_{2}^j(z) = \mathrm{Im}\,h^{j}(z), \quad j=1,2, \ z\in\SS. 
\end{equation}
Applying the mappings $\PP_{-}, \PP_{-}^{2}, \ldots, \PP_{-}^{n}$ to \eqref{formula-1}, where $n\ge 1$ is an even number, and using \eqref{eq-p1}, we obtain the following sequence of equations
\begin{equation}
\begin{aligned}\label{formula-2}
\tilde{w}_{2}^j(z) - \tilde{w}_{2}^j(\PP_{-}(z)) & = \mathrm{Im}\, h^j(\PP_{-}(z)), \\
\tilde{w}_{2}^j(\PP_{-}(z)) - \tilde{w}_{2}^j(\PP_{-}^{2}(z)) & = \mathrm{Im}\, h^j(\PP_{-}^{2}(z)), \\
& \vdots \\
\tilde{w}_{2}^j(\PP_{-}^{n-1}(z)) - \tilde{w}_{2}^j(\PP_{-}^{n}(z)) & = \mathrm{Im}\, h^j(\PP_{-}^{n}(z)).
\end{aligned}
\end{equation}
Adding the formulas \eqref{formula-2} side-by-side and using \eqref{eq-p2}, we obtain
\begin{equation*}
\begin{aligned}
\tilde{w}_{2}^j(z) & = \tilde{w}_{2}^j(\PP_{-}^{n}(z)) + \sum_{k=1}^{n}\mathrm{Im}\, h^j(\PP_{-}^{k}(z)) \\
& = \tilde{w}_{2}^j\left(z - \frac{2\pi n i}{1+a^2}\right) + \sum_{k=1}^{n}\mathrm{Im}\, h^j(\PP_{-}^{k}(z))
\end{aligned}
\end{equation*}
for $j=1,2$ and even $n\ge 1$. Passing in the above equation to the limit with $n\to \infty$ and using the assumption \eqref{obietnica}, we conclude that 
\begin{align*}
\tilde{w}_{2}^j(z) = \sum_{k=0}^{\infty}\mathrm{Im}\, h^j(\PP_{-}^{k}(z)), \quad j=1,2, \ z\in\SS.
\end{align*}
Recalling that $h^1(z) = h^2(z)$ for $z\in\SS$, we obtain
\begin{align*}
\tilde{w}_{2}^1(z) & = \sum_{k=0}^{\infty}\mathrm{Im}\, h^1(\PP_{-}^{k}(z)) = \sum_{k=0}^{\infty}\mathrm{Im}\, h^2(\PP_{-}^{k}(z)) = \tilde{w}_{2}^2(z), \quad z\in\SS.
\end{align*}
It remains to show that the real parts of the velocity fields $\tilde w^{j}$ coincide. Since the function $(\tilde{w}^{1}-\tilde{w}^{2})^{*}$ is holomorphic on the strip $\SS$ and its imaginary part vanishes, from the Cauchy-Riemann equation we obtain $$(\tilde{w}^{1}_{1})_{y}(z) = (\tilde{w}^{2}_{1})_{y}(z), \quad z\in\SS,$$ which together with \eqref{obietnica} gives
$\tilde{w}^{1}_{1} = \tilde{w}^{2}_{1}$ on $\SS$ as desired. Thus the proof of the proposition is completed. 
\end{proof}

\end{document}
